% =====================================================================
%  src/preamble.tex  —  导言区（SY05-24 Gadde《超共形指标讲义》）
%  由 Lead 维护；译者只读。缺宏时在文末「追加区」append \newcommand，
%  并在回报里列出新增宏名（不得改动上面的已有定义）。
% =====================================================================
\usepackage[a4paper,margin=2.6cm,top=3.0cm,bottom=2.8cm]{geometry}
\usepackage{amsmath,amssymb,amsthm,mathtools}
\usepackage{mathrsfs}
\usepackage{graphicx}
\usepackage{array,booktabs,longtable,multirow}
\usepackage{caption}
\usepackage{enumitem}
\usepackage{xcolor}
\usepackage{fancyhdr}

% ---------- 字体（Windows，显式路径） ----------
\setCJKmainfont[Path=C:/Windows/Fonts/,BoldFont=simhei.ttf,ItalicFont=simkai.ttf]{simsun.ttc}
\setCJKsansfont[Path=C:/Windows/Fonts/]{simhei.ttf}
\setCJKmonofont[Path=C:/Windows/Fonts/]{simfang.ttf}
\setCJKfamilyfont{zhkai}[Path=C:/Windows/Fonts/]{simkai.ttf}
\setCJKfamilyfont{zhhei}[Path=C:/Windows/Fonts/]{simhei.ttf}
\setCJKfamilyfont{zhfang}[Path=C:/Windows/Fonts/]{simfang.ttf}
\providecommand{\kaishu}{\CJKfamily{zhkai}}
\providecommand{\heiti}{\CJKfamily{zhhei}}
\providecommand{\fangsong}{\CJKfamily{zhfang}}

% ---------- 数学宏（块文件里只许用这些，禁止自造） ----------
\newcommand{\sN}{\mathcal{N}}          % 超对称数 \sN=1, \sN=2, \sN=4
\newcommand{\sI}{\mathcal{I}}          % 指标（index）
\newcommand{\sO}{\mathcal{O}}          % 算符
\newcommand{\sH}{\mathcal{H}}          % 希尔伯特空间
\newcommand{\sZ}{\mathcal{Z}}          % 配分函数
\newcommand{\sL}{\mathcal{L}}          % 拉氏量
\newcommand{\sQ}{\mathcal{Q}}          % 超荷
\newcommand{\sU}{\mathcal{U}}
\newcommand{\sV}{\mathcal{V}}
\DeclareMathOperator{\Tr}{Tr}
\DeclareMathOperator{\PE}{PE}
\DeclareMathOperator{\Sym}{Sym}
\DeclareMathOperator{\adj}{adj}
\DeclareMathOperator{\fund}{fund}
\DeclareMathOperator{\sgn}{sgn}
\DeclareMathOperator{\Li}{Li}
\DeclareMathOperator{\SU}{SU}
\DeclareMathOperator{\SO}{SO}
\DeclareMathOperator{\Sp}{Sp}
\DeclareMathOperator{\USp}{USp}
\DeclareMathOperator{\U}{U}
\DeclareMathOperator{\GL}{GL}
\DeclareMathOperator{\diag}{diag}
\newcommand{\R}{\mathbb{R}}
\newcommand{\C}{\mathbb{C}}
\newcommand{\Z}{\mathbb{Z}}
\newcommand{\Q}{\mathbb{Q}}
\newcommand{\bbZ}{Z}
\newcommand{\eps}{\varepsilon}
\newcommand{\dd}{\mathrm{d}}
\newcommand{\ee}{\mathrm{e}}
\newcommand{\ii}{\mathrm{i}}
\newcommand{\ol}[1]{\overline{#1}}
\newcommand{\br}[1]{\bar{#1}}
\newcommand{\dg}{\dagger}
\newcommand{\half}{\tfrac{1}{2}}
\newcommand{\third}{\tfrac{1}{3}}
\newcommand{\twothird}{\tfrac{2}{3}}
% 表格里常用的列类型
\newcolumntype{C}[1]{>{\centering\arraybackslash}p{#1}}

% ---------- 定理类环境（编号一律手工给：\begin{exercise}[2.1]） ----------
\theoremstyle{definition}
\newtheorem*{exercise}{习题}
\newtheorem*{definition}{定义}
\newtheorem*{example}{例}
\newtheorem*{theorem}{定理}
\newtheorem*{proposition}{命题}
\newtheorem*{lemma}{引理}
\newtheorem*{corollary}{推论}
\newtheorem*{remark}{注}
\newtheorem*{solution}{解}

% 让 \begin{exercise}[2.1] 印成「习题 2.1」（无括号）
\makeatletter
\def\thmhead@plain#1#2#3{%
  \thmname{#1}\thmnumber{\@ifnotempty{#2}{ #2}}%
  \thmnote{\@ifnotempty{#3}{\hspace{0.35em}{#3}}}}
\let\thmhead\thmhead@plain
\makeatother

% ---------- 版式 ----------
\setlength{\parindent}{2em}
\linespread{1.18}
\setlist{nosep,leftmargin=2.2em}
\renewcommand{\arraystretch}{1.15}

\pagestyle{fancy}
\fancyhf{}
\fancyhead[L]{\small\kaishu 超共形指标讲义}
\fancyhead[R]{\small\thepage}
\renewcommand{\headrulewidth}{0.4pt}
\fancypagestyle{plain}{\fancyhf{}\fancyfoot[C]{\small\thepage}\renewcommand{\headrulewidth}{0pt}}

\ctexset{
  section = {
    format = \Large\bfseries\heiti,
    beforeskip = 2.2ex plus 1ex minus .2ex,
    afterskip  = 1.4ex plus .2ex,
    name = {},
    number = \arabic{section},
  },
  subsection = {
    format = \large\bfseries\heiti,
    beforeskip = 1.8ex plus .6ex minus .2ex,
    afterskip  = 1.0ex plus .2ex,
  },
}

\usepackage[hidelinks,bookmarksnumbered=false,bookmarksopen=true,
            pdftitle={超共形指标讲义（中文译注本）},
            pdfauthor={Abhijit Gadde}]{hyperref}
\pdfstringdefDisableCommands{\def\ol#1{#1}}

% ===================== 追加区（译者可 append \newcommand） =====================
